\section{Comparison with chromatic homotopy theory}
\label{comp:section}

The algebraic construction and its obstruction are now complete. To
turn them into a spectrum, we use the large-prime comparison between
$E(2)$-local homotopy theory and the periodic derived category on
$\mathcal M$. Two identifications make that comparison useful here:
the localized sphere must correspond to the pure unit, and the
localized truncated Brown--Peterson spectrum must correspond to its
coefficient-chart sheaf. The first recovers ordinary homotopy groups
from algebraic sections. The second carries a proposed finite
construction into the algebraic setting, where the coefficient-change
argument will compare it with the obstruction on $\mathcal N$.

\subsection{The exact comparison}

Let $E=E(2)$ be the $p$-local Johnson--Wilson theory, using the
Hazewinkel coefficient chart from Proposition~\ref{geom:law-charts}.
In this section let $p>8$ and $D=2(p-1)$; in particular $D>12$.
This bound concerns the comparison. The operator construction imposes
the larger prime bounds already used in the algebraic argument.

\begin{theorem}[Large-prime algebraicity]
\label{comp:algebraicity}
For $p>8$ there is a triangulated equivalence
\[
 \mathcal F:h\Sp_E\longrightarrow \mathbf D(\mathcal M)
\]
whose sheaf homology is naturally identified, under descent and the
periodic regrading below, with the $E_*E$-comodule homology functor
$E_*(-)$. Its suspension structure
gives natural isomorphisms
\[
 \mathcal F(\Sigma^jZ)\cong(\mathcal F Z)[j]
 \qquad(j\in\Z).
\]
\end{theorem}

We use the construction in the November 14, 2018 version (v2) of
Pstr\k{a}gowski's preprint. Its Corollary~6.9 gives the range
$p>n^2+n+2$, hence $p>8$ at height two, by comparing homotopy
2-categories \cite[Corollary~6.9]{Pstragowski2021}. Two inputs to
that construction need explicit support here: the bound $6$ for Ext
between all graded comodules, and convergence of the topological and
algebraic potential-stage towers. The quotient-closure inference in
the proof of v2 Theorem~2.4 does not follow from the long exact Ext
sequence, and Proposition~4.4 defers the needed convergence. We first
identify the target, then prove these inputs for this fixed $E(2)$,
and finally follow the stage-zero homology and triangulated structures
through the same comparison.

Write $\mathcal D(E_*E)$ for the stable localization of differential
$E_*E$-comodules at homology isomorphisms, as in v2
Definitions~3.1--3.2 and Proposition~3.3
\cite{Pstragowski2021}. Its homotopy category is our twisted periodic
model by the following degree regrading. All coefficient and
cooperation degrees are multiples of $D$. For a differential comodule
$M$ in the homological convention, so that its differential lowers
total degree by one, define the weight-$s$ part of $C_j$ by
\[
 (C_j)_s=M_{j+sD}\qquad(j,s\in\Z).
\]
The graded comodule structure becomes the weight equivariance for flat
descent to $\mathcal M$. The definition gives
$(C_{j+D})_s=(C_j)_{s+1}$, hence $C_{j+D}=C_j(1)$, and the differential
gives $C_j\to C_{j-1}$. Conversely, write each total degree uniquely as
$n=r+sD$ with $0\leq r<D$ and set $M_n=(C_r)_s$. For $r>0$ the
differential is the component of $C_r\to C_{r-1}$. At $r=0$, the
identity $C_{-1}=C_{D-1}(-1)$ identifies its target with
$(C_{D-1})_{s-1}$, the component of total degree $n-1$; this is the
required twist at the residue boundary. The coefficient action and
coaction are recovered from the weight-graded ones. Degree-zero maps
restrict to families of sheaf maps commuting with differentials and
periodicity, and the inverse formula recovers every such map uniquely.
The two constructions identify homology in every degree and weight,
so they identify quasi-isomorphisms. They also identify shifts and
mapping cones: the suspension convention is
$(C[h])_j=C_{j-h}$ with differential multiplied by $(-1)^h$.
In particular $C[D]\cong C(-1)$ with its differential sign, since
$D$ is even. Thus regrading gives an exact equivalence
$h\mathcal D(E_*E)\simeq\mathbf D(\mathcal M)$. It applies to all
differential comodules and periodic complexes, including arbitrary
unbounded and non-finite-type objects.

\begin{lemma}[The height-two comparison inputs]
\label{rew:comp:height-two-inputs}
Let $\mathscr C_E=\operatorname{Comod}_{E_*E}$ be the ordinary abelian
category of all graded comodules. Then
\[
 \operatorname{Ext}_{\mathscr C_E}^{s}(M,N)=0
 \qquad(M,N\in\mathscr C_E,\ s>6).
\]
This includes every internal grading shift. Let $\mathrm{Syn}_E$ be
the Grothendieck prestable infinity-category of spherical hypercomplete
sheaves of spaces on the site $\Sp_E^{\mathrm{fp}}$ of finite spectra
$P$ with $E_*P$ projective over $E_*$, with covers the single maps
surjective on $E$-homology. Its objects are connective synthetic
spectra, and each satisfies
\[
 X\simeq\lim_{i\geq0}X_{\leq i}\qquad(X\in\mathrm{Syn}_E).
\]
The ordinary unbounded derived infinity-category
$\mathcal D_{\mathrm{ord}}(\mathscr C_E)$ is left complete as well.
For the topological potential stages of the unit shift algebra on this
same $\mathrm{Syn}_E$, and the algebraic potential stages of the
periodicity algebra on $\mathcal D_{\mathrm{ord}}(\mathscr C_E)_{\geq0}$,
the natural functors are equivalences
\[
 \mathcal M_\infty^{\mathrm{top}}\simeq
       \lim_\ell\mathcal M_\ell^{\mathrm{top}},\qquad
 \mathcal M_\infty^{\mathrm{alg}}\simeq
       \lim_\ell\mathcal M_\ell^{\mathrm{alg}}.
\]
These are the two convergence inputs $(\spadesuit5)$ of v2
Propositions~4.4 and~4.8.
\end{lemma}

\begin{proof}
Put
\[
 \begin{gathered}
 B=E_*=\Z_{(p)}[v_1^{\mathrm{Haz}},(v_2^{\mathrm{Haz}})^{\pm1}],
 \\
 I_0=(0),\quad I_1=(p),\quad
 I_2=(p,v_1^{\mathrm{Haz}}),\quad Q_h=B/I_h\quad(0\leq h\leq2).
 \end{gathered}
\]
The ideals $I_h$ are invariant under the chart Hopf algebroid, by
\eqref{geom:invariant-ideals} and the triangular coordinate change.
Because the coefficient and cooperation degrees lie in $D\Z$, the
degree phases split $\mathscr C_E$ into a product of $D$ abelian
categories. After shifting a phase and reindexing its degrees as
weights, flat descent on the chart of Proposition~\ref{geom:law-charts}
identifies it with $\operatorname{QCoh}(\mathcal M)$. This is abelian
descent, independent of the differential regrading and of topological
algebraicity.

\paragraph{Supported unit-source bounds.}
For a comodule $V$ annihilated by $I_h$, the supported
unit-source bounds are
\[
 \operatorname{Ext}_{\mathscr C_E}^{a}(B,V)=0\qquad(a>6-h),
\]
also after arbitrary internal shifts. For $h=0$ this is the bound on
$\mathcal M$ in Proposition~\ref{geom:cohomological-bounds}. For $h=1$,
the closed-base-compatible coinduced resolution in
Lemma~\ref{geom:sections-resolution} identifies the sections
calculation with the characteristic-$p$ one, where that proposition
gives the bound $5$.

For $h=2$, regard the corresponding sheaf as a sheaf on
$\mathcal M_0$ and extend it to $\mathcal M_k$, writing $V_k$ for the
result. The congruence proved in Remark~\ref{comp:forms} identifies
$v_1^{\mathrm{Haz}}$ modulo $p$ with the Cartier section $t$, so
$tV_k=0$. The central action is $\mathcal O$-linear. Consequently
both $r_*V_k=\ker(c-1)$ and $R^1r_*V_k=\coker(c-1)$ are still killed
by $t$, and the higher derived invariants vanish, by
Proposition~\ref{geom:cohomological-bounds}. Since $r^*$ is exact,
$r_*$ preserves injectives, and the derived-composite spectral
sequence is
\[
 E_2^{a,b}=H^a(\mathcal N,R^b r_*V_k)
 \Longrightarrow H^{a+b}(\mathcal M_k,V_k).
\]
There are only rows $b=0,1$. The closed-base resolution and
Lemma~\ref{geom:group-cohomology} apply to both rows on
$\mathcal N/(t)$ and give vanishing for $a>3$, for arbitrary sheaves.
Thus the total degree is at most~$4$. Finite flat extension
$\F_p\to k$ commutes with the coinduced sections resolution and
detects vanishing. The closed-base compatibility of that resolution
then returns the bound from $\mathcal M_0$ to $\mathcal M$. These three
bounds are uniform under weight
twists. They are bounds with unit source; the following argument is
needed to obtain the all-pairs assertion.

\paragraph{From supported units to all pairs.}
The cooperation Hopf algebroid is flat and Adams. Indeed the empty
regular sequence in Hovey's Proposition~1.4.8 gives topological
flatness of $\BP$; Theorem~1.4.9 carries it to its Landweber-exact
module $E$, and Lemma~1.4.6 uses the commutative cooperation algebra
to give the Adams property \cite{Hovey2004}. Propositions~1.4.4
and~1.4.1 there imply that $\mathscr C_E$ is Grothendieck, that its
dualizable comodules generate, and that every comodule finitely
generated over $B$ is a quotient of a dualizable one. By
Proposition~1.3.4 there, dualizability means that the underlying
$B$-module is finitely generated projective.

The ring $B$ is noetherian, and the regular sequence defining $I_h$
gives $Q_h$ a graded $B$-module Koszul resolution of length $h$.
Starting with $K_0=Q_h$, choose dualizable epimorphisms
$P_j\to K_j$ and put $K_{j+1}=\ker(P_j\to K_j)$ for $0\leq j<h$.
Noetherianity keeps these kernels finitely generated over $B$.
For every graded $B$-module $W$, dimension shifting gives
\[
 \operatorname{Ext}^1_B(K_h,W)
 \cong\operatorname{Ext}^{h+1}_B(Q_h,W)=0.
\]
Thus $K_h$ is projective over $B$ and is dualizable. For $h=0$ the
object $Q_0=B$ already is dualizable. This constructs in every case
a length-$h$ resolution $P^\bullet\to Q_h$ by dualizable comodules,
in ordinary cohomological degrees $-h,\ldots,0$.

For an arbitrary comodule $V$, the bounded complex
$(P^\bullet)^\vee\otimes_B V$ computes the internal derived Hom from
$Q_h$ in $\mathcal D_{\mathrm{ord}}(\mathscr C_E)$. To see this within
that category, each dualizable term has evaluation and coevaluation,
and its tensor functor is exact; bounded totalization therefore
preserves quasi-isomorphisms and the duality descends to the ordinary
derived category. The underlying complex is the graded
$\operatorname{Hom}_B(P^\bullet,V)$. Its cohomology in degree $b$ is
$\operatorname{Ext}^b_B(Q_h,V)$, including its internal grading, and
vanishes outside $0\leq b\leq h$. It is annihilated by $I_h$: the
$B$-action on this Ext can be taken through its source $Q_h$, or,
equivalently, multiplication by an element of $I_h$ lifts to a
null-homotopic map on a projective resolution. Since $I_h$ is
invariant, these cohomology objects are supported comodules. The
bounded hypercohomology spectral sequence is consequently
\[
 E_2^{a,b}=\operatorname{Ext}_{\mathscr C_E}^{a}
   \bigl(B,H^b((P^\bullet)^\vee\otimes_B V)\bigr)
 \Longrightarrow \operatorname{Ext}_{\mathscr C_E}^{a+b}(Q_h,V).
\]
Its entries have $a\leq6-h$ by the supported unit bounds and
$0\leq b\leq h$ by the resolution. Hence
$\operatorname{Ext}^{>6}_{\mathscr C_E}(Q_h,V)=0$ for arbitrary $V$,
uniformly in internal shifts. This calculation uses a bounded
dualizable comodule resolution in the ordinary derived category; it
requires no comparison with a different derived category of sheaves.

By the Landweber filtration theorem
\cite[Theorem~5.7]{HS2005}, every comodule finitely presented over $B$
has a finite filtration by internal shifts of $Q_0,Q_1,Q_2$. The long
exact Ext sequence therefore gives the same bound $6$ for every such
source and every target. Let $\mathscr U$ be this class of finite
sources. It is quotient closed: a quotient remains finitely generated
over the noetherian ring $B$, hence finitely presented, and comodule
quotients have the underlying module quotients. It also detects every
nonzero comodule, because the primitive theorem
\cite[Theorem~5.1]{HS2005} supplies a nonzero map from an internal shift
of $B$ to each one.

We use the partial-section argument underlying Franke's
injective-truncation criterion
\cite[Section~3.4.3, Lemma~2]{Franke1996}. In a Grothendieck abelian
category, let a quotient-closed class $\mathscr U$ detect nonzero
objects. If $K$ is injective and $\pi:K\to L$ induces surjections
$\operatorname{Hom}(U,K)\to\operatorname{Hom}(U,L)$ for every
$U\in\mathscr U$, then $\pi$ admits a section. Indeed partial sections
on subobjects of $L$ form a set, by well-poweredness and local
smallness, and exact filtered unions give upper bounds for chains.
Choose a maximal partial section with domain $L_0$. If $L_0\ne L$,
choose a nonzero map from an object of $\mathscr U$ to $L/L_0$.
Its image $V$ still belongs to $\mathscr U$ by quotient closure; let
$L_1$ be its preimage in $L$. Injectivity extends the partial section
to a map $\widetilde s:L_1\to K$. The discrepancy
$\pi\widetilde s-(L_1\hookrightarrow L)$ vanishes on $L_0$, so it
factors through a map $V\to L$. Lift that map to $K$ by the Hom
surjection and subtract its composite with $L_1\to V$. This extends
the section to $L_1$, contradicting maximality.

For any target $N$, take an injective resolution
$0\to N\to I^0\to I^1\to\cdots$ and set
$Z^7=\ker(I^7\to I^8)$. Dimension shifting and the finite-source
bound show that $I^6\to Z^7$ induces surjections on
$\operatorname{Hom}(U,-)$ for every $U\in\mathscr U$: their cokernels
are $\operatorname{Ext}^7(U,N)=0$. The criterion splits this
epimorphism. Its kernel is an injective direct summand of $I^6$, so
the resolution truncates at degree $6$. Thus every target has
injective dimension at most $6$, proving the stated Ext bound for
all sources and targets.

\paragraph{Postnikov completeness.}
It remains to prove the two convergence assertions. We apply
Pstr\k{a}gowski--VanKoughnett's general synthetic construction and
the argument of their Theorem~7.4, with locators in arXiv v4
\cite{PV2022}. The theory $E(2)$ is Adams-type by their Example~6.2.
The site in the statement consists of the unlocalized finite spectra
$P$ and its stated cover maps. Taking all its objects and their
identity covers in their Lemma~7.2 gives the colimit-generating family
\[
 G_P=\nu(L_EP)\quad\text{in }\mathrm{Syn}_E.
\]
For $X\in\mathrm{Syn}_E$, put
$X[i]=\operatorname{fib}(X_{\leq i}\to X_{\leq i-1})
\simeq\Sigma^i\pi_iX$. This layer is canonically a module over the
unit truncation $\mathbf 1_{\leq0}$. Their Theorem~3.11 together with
Proposition~6.4 identifies that module category with
$\mathcal D_{\mathrm{ord}}(\mathscr C_E)_{\geq0}$, and identifies the
reduction of $G_P$ with $E_*P$. By the module adjunction, for $j\geq0$
we therefore have
\[
 \pi_j\operatorname{Map}_{\mathrm{Syn}_E}(G_P,X[i])
 \cong\operatorname{Ext}_{\mathscr C_E}^{i-j}(E_*P,\pi_iX).
\]
This vanishes when $i>6+j$, uniformly in $P$ and in the arbitrary
layer comodule $\pi_iX$. The fiber sequences of the Postnikov tower
therefore stabilize these mapping-space homotopy groups.

For precision, the completeness criterion in the proof of their
Theorem~7.4 asks that, for every $m\geq0$, the projection
$\lim_iX_{\leq i}\to X_{\leq k}$ become an $m$-equivalence for large
$k$, meaning an equivalence after $m$-truncation. It is enough that
the induced maps on mapping spaces out of the colimit generators be
$(m+1)$-equivalences.
Representable mapping spaces preserve limits, so the preceding
stabilization and the Milnor exact sequence give this condition.
Comparing with $X\to X_{\leq k}$ and using separatedness of
$\mathrm{Syn}_E$ proves $X\simeq\lim_iX_{\leq i}$. The literal
Theorem~7.4 is stated for Morava $E$-theory; this is its proof applied
on the stated $E(2)$ site, using Lemma~7.2 and the local uniform Ext
bound in place of its Morava-specific finite-source input.

\paragraph{Potential-stage convergence.}
This objectwise Postnikov completeness supplies the separate formal
tower step. Proposition~3.2 of the same source recovers unit-module
categories as the limit of the categories of modules over the unit
truncations; Definition~4.1 and Remark~4.4 restrict this equivalence
to the potential stages. Proposition~6.4 identifies the periodic
unit modules in this synthetic category with $\Sp_E$. On the
algebraic side the finite homological dimension just proved gives
left completeness of $\mathcal D_{\mathrm{ord}}(\mathscr C_E)$ by the
argument recorded in Remark~7.5. The same formal module and
potential-stage step applies to the algebraic periodicity algebra of
v2 Proposition~4.8. This proves both displayed tower equivalences in
the exact categories used by the comparison.
\end{proof}

\begin{proof}[Proof of Theorem~\ref{comp:algebraicity}]
We now follow the comparison in Pstr\k{a}gowski v2
\cite{Pstragowski2021}. Propositions~4.4 and~4.8, item
$(\spadesuit6)$, identify stage-zero reduction with $E$-homology on
the topological side and differential-comodule homology on the
algebraic side. Lemma~5.3 identifies the bottom functor on finitely
projective comodules with their inclusion into the standard comodule
heart, and its cocontinuous extension gives the heart identification.
Lemmas~5.6 and~5.8 identify the heart and the homotopy objects used by
the stages. Corollary~6.5 uses the preferred unit/free-algebra map and
commutes with reduction to stage zero; Corollary~6.7 restricts that
comparison to potential stages.

Lemma~\ref{rew:comp:height-two-inputs} supplies the dimension $d=6$
and the two convergence inputs used by v2 Theorem~4.10. Its
stabilization stage is $\ell=d+k-1$ at homotopy level $k$, hence
$\ell=7$ at level two. Theorem~6.8, in the stated Corollary~6.9 range
$p>8$, then gives the zigzag of homotopy 2-category equivalences
\[
 h_2\Sp_E\xrightarrow{\sim}h_2\mathcal M_7^{\mathrm{top}}
 \simeq h_2\mathcal M_7^{\mathrm{alg}}
 \xleftarrow{\sim}h_2\mathcal D(E_*E).
\]
Choose an inverse to the final leg. Compose from left to right and
then apply the regrading above on homotopy categories. Call the
resulting equivalence $\mathcal F$. The stage-zero compatibility
throughout the restriction and assembly gives its natural identification of
sheaf homology with $E_*(-)$.

The recovery of triangles uses Appendix~B of v2, with one local
correction: the proof of Lemma~B.3 overstates the local connectivity
of the projection from a stable infinity-category to its homotopy
2-category. Its needed uniqueness follows directly from the cone
mapping space. Write $\mathcal H=h_2\mathscr T$ for a stable
infinity-category $\mathscr T$, and fix a span
$\sigma=(B\leftarrow A\to C)$. The space of cones with vertex $W$ is
\[
 \operatorname{Cone}_{\sigma}(W)\simeq
 \operatorname{Map}_{\mathcal H}(B,W)
 \mathop{\times}^{h}_{\operatorname{Map}_{\mathcal H}(A,W)}
 \operatorname{Map}_{\mathcal H}(C,W),
\]
where the homotopy pullback retains the cone's 2-cell. If $q_c$ is a
cone with vertex $V$ and $q$ is any target cone with vertex $W$,
Lemma~A.5 identifies the mapping space of cones as
\[
 \operatorname{Map}_{\mathcal H_{\sigma/}}(q_c,q)
 \simeq\operatorname{hofib}_{q}\bigl(
 \operatorname{Map}_{\mathcal H}(V,W)
 \longrightarrow\operatorname{Cone}_{\sigma}(W)\bigr).
\]
The arrow postcomposes the entire cone $q_c$ with a vertex map.
When $q_c$ is 1-cocartesian, Definition~B.1 says that the displayed
comparison is 1-connected. Its fiber over the arbitrary cone $q$ is
therefore nonempty and connected, so it has exactly one $\pi_0$ class.
Existence follows by lifting the span to $\mathscr T$, forming a
pushout, and applying Lemma~B.2 to its image. This proves the
existence and full uniqueness needed in Lemma~B.3. Corollary~B.4 then
identifies any two such cones. The arbitrary-target form also gives
suspension maps and their naturality: for $f:X\to Y$, precompose the
chosen cone under $0\leftarrow Y\to0$ to obtain a cone under
$0\leftarrow X\to0$ with vertex $\Sigma Y$, and apply that uniqueness.
The two-dimensional lifting in Theorem~B.6, together with
Lemma~B.2 and Corollary~B.4, consequently recovers distinguished triangles,
their boundary classes, and the natural suspension comparison;
Corollary~B.7 makes the induced equivalence triangulated.

Apply this recovery to the chosen equivalence of homotopy
2-categories and then use the exact regrading. It equips the same
$\mathcal F$ with its triangulated structure. Iterating its natural
suspension isomorphism and its inverse gives the displayed
isomorphisms for every integer $j$.
\end{proof}

Franke introduced the differential-object model and realization
construction \cite[Section~2.2 and Theorem~10 of Section~3.5]{Franke1996}.
Patchkoria identified a gap in the general equivalence argument
\cite[Remark~6.3.1]{Patchkoria2011}. The comparison just justified is
Pstr\k{a}gowski's later Goerss--Hopkins construction; these results
do not identify it with Franke's particular functor.
Patchkoria--Pstr\k{a}gowski give a homology-theoretic formulation of
the comparison \cite[Theorems~7.56 and~8.13]{PP2023}, and Barkan proves
a tensor-triangulated strengthening in its stated prime range
\cite[Theorem~1.1 and Corollary~3.42]{Barkan2023}. We use the chosen
exact comparison, its homology identification, and distinguished
triangles and retracts. The inverse limits, completions, and gluing
used to construct the algebraic object were formed in the algebraic
homotopy theory before realization; no monoidal or stable
infinity-category comparison between the topological and algebraic
theories is used.

\subsection{The Brown--Peterson quotient and its comodule}

Recall the rings $P=\BP_*$ and $B_2=\Z_{(p)}[v_1,v_2]$, the chart
$f:\mathcal X\to\mathcal M$, and the sheaves
\[
 T_{\Z}=f_*\mathcal O_{\mathcal X},\qquad
 T_{\Z_p}=f_*\bigl(\mathcal O_{\mathcal X}
                  \otimes_{\Z_{(p)}}\Z_p\bigr).
\]
The comparison theorem identifies homology, but this alone does not
identify the image of a chosen spectrum with a pure sheaf. We first
calculate the comodule of a standard Brown--Peterson quotient; the
sparsity argument in the next subsection will then identify the
object itself.

We use the regular-quotient form $b$ of $\BP\langle2\rangle$.
The existence of an $\mathbb E_4$ multiplication on $\BP$
\cite[Theorem~1.1 and Corollary~1.2]{BM2013} gives a monoidal module
category in which to form the cofibers of the scalar maps $v_i$, $i>2$.
Taking these regular quotients
successively and then their sequential colimit gives a $\BP$-module
with coefficients $B_2$. Indeed the scalar action on the homotopy of a
module is multiplication by the corresponding coefficient: homotopy
classes are maps out of shifts of the monoidal unit, and tensor
naturality identifies the scalar action with composition at that unit.

We explain why the enhanced multiplication causes no change in the
cooperation formulas used below. An $\mathbb E_4$ ring with underlying
spectrum $\BP$ has even homotopy and admits a complex orientation.
The orientation determines a unital multiplicative map from
$MU_{(p)}$ in the homotopy category: its Thom classes multiply under
Whitney sum. Composing with the multiplicative inclusion of the Quillen
summand gives a homotopy ring map from the usual $\BP$ to this
enhancement; see \cite[Theorem~4.1.12 and Lemma~4.1.13]{Ravenel1986}.
Both $MU_{(p)}$ and $MU_{(p)}\wedge MU_{(p)}$ have filtrations by
finite even-cell skeleta. Maps from their suspended finite skeleta to
the even target vanish by cellular induction. The corresponding Milnor
$\lim^1$ terms therefore vanish, so the compatible Thom maps and their
multiplicativity identities pass to the whole spectra without phantom
ambiguities.
This is the even-cell case of the no-even-phantoms criterion in
\cite[Lemma~2.11, Corollary~2.15, and Proposition~2.18]{HoveyStrickland1999}.

This map is an equivalence. It induces a unit on degree-zero mod-$p$
cohomology, since the unit of $\pi_0\cong\Z_{(p)}$ is not
$p$-divisible. The Steenrod module $H^*(\BP;\F_p)$ is cyclic on this
class, and source and target have the same finite dimensions in each
degree. The induced map is therefore an isomorphism on mod-$p$
cohomology. Its cofiber is connective and of finite type over
$\Z_{(p)}$; lowest-degree Hurewicz and Nakayama show that this cofiber
vanishes. We may accordingly choose the usual coefficient generators
and cooperation formulas for the multiplication acting on these modules.
Naturality of multiplication and of the unit insertions identifies the
units and coproduct on cooperations under this homotopy ring equivalence.
All module constructions remain in the module category of the chosen
$\mathbb E_4$ enhancement; this argument requires no refinement of the
comparison map to a map of structured rings.

\begin{proposition}[The quotient comodule]
\label{comp:bp-comodule}
For the regular quotient $b$ just constructed,
\[
 \BP_*b\cong\Gamma_P\otimes_{\eta_R,P}B_2.
\]
Under Landweber base change to $E$ and descent to $\mathcal M$, this
comodule is the pure sheaf $T_{\Z}$. The corresponding comodule for
$b_p^\wedge$ is $T_{\Z_p}$. Both sheaves are $p$-monic, and their
reductions modulo $p$ equal $T_0$.
\end{proposition}

\begin{proof}
There is a relative smash identification
\[
 \BP\wedge b\simeq
 (\BP\wedge\BP)\otimes_{\BP}b.
\]
As a right $\BP$-module, $\BP\wedge\BP$ is a wedge of free shifts:
lift a right-module basis of its homotopy to obtain the equivalence.
Thus the flat K\"unneth calculation gives the asserted tensor product.
Under that identification, insertion of units, which defines the
coaction, agrees with the cooperation coproduct. Only the homotopy
cooperation formulas, not a further multiplicative enhancement, enter
this calculation.

Landweber exactness now gives, on the Hazewinkel chart,
\[
 A_{\rm Haz}\otimes_{P,\eta_L}\Gamma_P
                  \otimes_{\eta_R,P}B_2.
\]
These are the functions on strict arrows with the chosen chart law at
the left and a $\mathcal X$-law at the right, with grading frames
identified. Descent is composition of arrows. This is precisely the
pullback of $f_*\mathcal O_{\mathcal X}$. All degrees lie in
$D\Z$, so its regraded sheaf is pure. Since $f$ is affine and
$\mathcal O_{\mathcal X}$ is $p$-monic, so is $T_{\Z}$.

The same calculation also identifies the completed generator. The
completion $b_p^\wedge$ is the homotopy limit of $b/p^a$.
The tower is surjective on homotopy, so its coefficients are
$B_2\otimes_{\Z_{(p)}}\Z_p$, with ordinary completion in each
graded degree. The limit is complete by its Postnikov description and
has the same multiplication-by-$p$ cofiber as $b$. It remains a
$\BP$-module, and the same free-right-module calculation gives
$T_{\Z_p}$. Its $p$-monicity and reduction follow on the affine chart
and then by affine pushforward.
\end{proof}

We have now matched the quotient coefficients and coaction to the
chart. To ensure that this particular regular quotient is a standard
truncated Brown--Peterson spectrum, we check its Steenrod module.

\begin{proposition}[Standard mod-$p$ cohomology]
\label{comp:standard-cohomology}
At an odd prime, the quotient $b$ has the standard Steenrod module
\[
 H^*(b;\F_p)\cong
 \mathcal A//E(Q_0,Q_1,Q_2),
\]
where $Q_i$ are the Milnor primitives and $E(Q_0,Q_1,Q_2)$ is their
exterior subalgebra. In particular its completion is a standard form
of $\BP\langle2\rangle_p^\wedge$.
\end{proposition}

The displayed Steenrod module is classical
\cite[Proposition~1.7]{Wilson1975}; we verify it for the module quotient
chosen here.

\begin{proof}
Connective $\BP$-module truncation and reduction give a map
$b\to H\F_p$. For $N=b,H\F_p$, consider the connective K\"unneth
spectral sequence for
\[
 H\F_p\otimes_{\BP}(\BP\wedge N),
\]
where $\BP$ acts on the first factor of $\BP\wedge N$.
The free-right-module calculation in
Proposition~\ref{comp:bp-comodule} applies to both choices of $N$.
The map on Tor terms is therefore
\[
 \F_p[t_1,t_2,\ldots]\otimes
       \bigwedge(\epsilon_i:i>2)
 \ \longrightarrow\
 \F_p[t_1,t_2,\ldots]\otimes
       \bigwedge(\epsilon_i:i\geq0),
 \qquad |\epsilon_i|=2p^i-1,
\]
and is the inclusion. To compute it, use the regular Koszul resolutions
of the right-base-change quotients of $\Gamma_P$, on
$\eta_Rv_i$ for $i>2$, and on $p,\eta_Rv_i$ for $i\geq1$,
respectively. These resolutions are free over the left base. After
tensoring with the left $\F_p$, every differential vanishes.

The target spectral sequence collapses by the dimensions in Milnor's
dual Steenrod algebra calculation, recalled in
\cite[Theorem~3.1.1(b)]{Ravenel1986}.
The injective map and naturality then force the source to collapse
page by page. Connectivity gives convergence, and the resulting map
on mod-$p$ homology is injective. Dually $\mathcal A$ surjects onto
$H^*(b;\F_p)$. There are no classes in the odd degrees of
$Q_0,Q_1,Q_2$, so these primitives kill the degree-zero generator.
The resulting quotient map from
$\mathcal A//E(Q_0,Q_1,Q_2)$ is an isomorphism by the same dimension
calculation.
\end{proof}

\begin{remark}[Choice of generators and forms]
\label{comp:forms}
The quotient comparison can equally be made with the Hazewinkel
chart. Its characteristic-$p$ reduction agrees with the Araki chart
used here. The congruence between the Araki and Hazewinkel generators
modulo $p$ is classical \cite[Theorem~A2.2.3]{Ravenel1986}; we give the
calculation in our notation. If $l_i$ are the logarithm coefficients, subtraction
of the two logarithm recursions gives
\[
 v_n-v_n^{\rm Haz}
 =-p^{p^n}l_n+
   \sum_{0<i<n}l_i\bigl((v_{n-i}^{\rm Haz})^{p^i}
                              -v_{n-i}^{p^i}\bigr).
\]
The facts that $p^il_i\in P$ and that raising to $p^i$ improves a
congruence modulo $p$ to one modulo $p^{i+1}$ prove inductively that
$v_n\equiv v_n^{\rm Haz}\pmod p$. Thus the subsequent reduction
argument does not depend on this coordinate choice.

The quotient $b$ is connective and each $\pi_i b$ is finitely generated
over $\Z_{(p)}$. Together with the Steenrod-module calculation above,
the completed uniqueness theorem of Angeltveit--Lind \cite{AL2017},
with Lee's corrected proof, identifies $b_p^\wedge$ with the standard
underlying spectrum $\BP\langle2\rangle_p^\wedge$
\cite[Theorem~1.1 and Remark~1.4]{Lee2026}; see also
\cite[Note~2]{LP2026}. For independence of the quotient generators
already before completion, see \cite[Corollary~1.9]{Lee2026}. That
statement concerns underlying spectra; the $\BP$-module structures
may differ. For a refutation of the universal thick generation
assertion, the specified standard quotient form already suffices.
\end{remark}

\subsection{Identifying the pure representatives}

We have identified the homology of the comparison objects. The
remaining step is to rule out nontrivial attachments between copies
of that homology in different periods. The cohomological bounds from
Section~\ref{geom:section} do this: the next lemma makes precise how a
gap in homology forces a periodic complex to be pure.

\begin{lemma}[Sparse representatives]
\label{comp:sparse}
In a twisted periodic derived category of period $D$, suppose a complex
has homology only in one residue modulo $D$, with sheaf $W$ in a
chosen representative of that residue. It is isomorphic to the
corresponding pure object in either of the following cases:
\begin{enumerate}
\item $W$ has projective dimension at most $a$ in the Ext sense and
      $D>a+1$;
\item $W$ has injective dimension at most $b$ and $D>b+1$.
\end{enumerate}
The assertion applies to every shift of the complex.
\end{lemma}

\begin{proof}
Shift the chosen residue to degree zero. For the first assertion, take
a fibrant representative $C$ with injective terms and write $B_j,Z_j$
for its boundaries and cycles. Lifting the isomorphism
$W\to H_0C$ to $Z_0C$ is obstructed by
$\Ext^1(W,B_0C)$. In the intervening positive degrees the homology
vanishes, so the short exact sequences
\[
 0\longrightarrow B_jC\longrightarrow C_j
       \longrightarrow B_{j-1}C\longrightarrow0
\]
dimension-shift this obstruction through injective terms. Before
reaching the next nonzero homology residue it becomes an Ext group
of degree greater than $a$, and hence vanishes. A lift gives a map
from the pure $W$ into $C$. Repeat it by the period-twist
identifications. It induces the given isomorphism on every homology
sheaf and is a quasi-isomorphism.

For the second assertion take an injective resolution of $W$ in
degrees $0,-1,\ldots,-b$, and periodize it. Map $C$ into this
resolution. The isomorphism on $Z_0C/B_0C$ extends to $C_0/B_0C$
with values in the first injective. Its composite with the next
differential vanishes on cycles, so it prescribes a map on $B_{-1}C$;
extend that map to $C_{-1}$ by injectivity. Continue in this way.
At every intermediate step, vanishing of homology identifies the
cycles with the boundaries on which the chain-map equation has
already been prescribed. The finite resolution ends before the next
homology residue, so its periodically repeated copies do not
interfere. The resulting map is again a quasi-isomorphism. The same
arguments begin in any chosen residue, proving the assertion for all
shifts.
\end{proof}

\begin{proposition}[The comparison generators]
For every integer $j$,\label{comp:generators}
\[
 \begin{split}
 \mathcal F(\Sigma^jL_ES^0)&\cong\mathcal O[j],\\
 \mathcal F(\Sigma^jL_Eb)&\cong T_{\Z}[j],\\
 \mathcal F(\Sigma^jL_Eb_p^\wedge)&\cong T_{\Z_p}[j].
 \end{split}
\]
\end{proposition}

\begin{proof}
The sphere has pure comodule homology $\mathcal O$, and
Proposition~\ref{geom:cohomological-bounds} gives
$\Ext^{>6}(\mathcal O,-)=0$. Thus the first part of
Lemma~\ref{comp:sparse} identifies its image, since $D>12>6+1$.

For the chart objects we claim injective dimension at most $5$.
Resolve $\mathcal O_{\mathcal X}$ by its sheaf \v{C}ech resolution for
the three principal opens, of length $2$. On each quotient-affine
chart the localized regular polynomial ring has global dimension
at most $3$. Hence each term has injective dimension at most $3$
after pushforward. This assertion applies to arbitrary graded
modules: a graded free resolution is also an ungraded free
resolution, and in its Hom complex the projection of every
homogeneous input to the output of matching degree splits off the
grading-preserving Hom complex. Its Ext groups are therefore
summands of the ordinary Ext groups. Flat-affine adjunction for $f$
preserves the injective-dimension bound. Combining the length-$2$
resolution with the bound $3$ gives the bound $5$ for $T_{\Z}$.

The same argument applies after tensoring the chart sheaf with
$\Z_p$: global dimension bounds hold for arbitrary modules over
the localized regular rings, so no finite-generation assertion over
$\Z_{(p)}$ is needed. It gives injective dimension at most $5$ for
$T_{\Z_p}$. Proposition~\ref{comp:bp-comodule}, homology
compatibility, and the second part of Lemma~\ref{comp:sparse}
now give the remaining identifications, since $D>5+1$. The bounds
are unchanged by the weight twists, so the sparsity proof applies in
every residue. Equivalently, once the degree-zero objects have been
identified, the suspension structure of $\mathcal F$ gives the
displayed isomorphisms for every $j$.
\end{proof}

Proposition~\ref{comp:generators} completes the two identifications
needed for transfer. For every realized algebraic object, maps from
the pure unit compute ordinary homotopy groups. Moreover, exactness of
$\mathcal F$ sends every finite construction from $L_Eb_p^\wedge$
to a finite construction from $T_{\Z_p}$: it carries shifts,
distinguished triangles, and retracts to the corresponding operations. The transfer
argument will use precisely these two consequences.

The main counterexample proof continues in
Section~\ref{transfer:section}. The following subsection gives an
independent categorical extension.

\subsection{A transport principle without exactness}
\label{rew:comp:weaker-transport}

There is a useful weaker principle for comparisons that are known
only as equivalences of underlying categories. In that setting the
functor need not identify suspension, so all shifts must appear in
the membership condition. This principle is independent of the exact
comparison chosen above.

\begin{lemma}[Transport of thick membership]
\label{comp:thick-transport}
Let $F:\mathcal C\to\mathcal D$ be an equivalence of the underlying
categories of triangulated categories, and let $\mathcal D_0$ be a
thick subcategory of $\mathcal D$. Then
\[
 \{X\in\mathcal C:
       F(\Sigma^jX)\in\mathcal D_0\text{ for every }j\in\Z\}
\]
is a thick subcategory of $\mathcal C$.
\end{lemma}

\begin{proof}
Equivalences preserve zero objects, finite biproducts, and retracts.
Closure under shifts is built into the definition. To prove closure
under triangles, consider
\[
 A\longrightarrow B\longrightarrow C\longrightarrow\Sigma A
\]
and first suppose only that $F(B),F(\Sigma A)\in\mathcal D_0$.
Complete the map $F(B)\to F(C)$ to a triangle
\[
 F(B)\longrightarrow F(C)\xrightarrow{s}D_1\longrightarrow F(B)[1].
\]
The map $s$ kills $F(B)\to F(C)$. Full faithfulness and essential
surjectivity transport the exactness of represented sequences in
$\mathcal C$, so $s$ factors through the image of the next map:
\[
 s=hF(C\to\Sigma A),\qquad h:F(\Sigma A)\longrightarrow D_1.
\]
Form a weak homotopy pullback $P$ of $s$ along $h$. Its triangle is
\[
 F(B)\longrightarrow P\longrightarrow F(\Sigma A)
     \longrightarrow F(B)[1],
\]
so $P\in\mathcal D_0$. The pair consisting of the identity of $F(C)$
and $F(C\to\Sigma A)$ has equal composites to $D_1$. The exact
sequence defining the weak pullback therefore gives a map
$F(C)\to P$ whose composite with $P\to F(C)$ is the identity.
Thus $F(C)$ is a retract of $P$ and belongs to $\mathcal D_0$.
Weak pullbacks here can be constructed from fiber triangles and the
octahedral axiom.

Now suppose $A$ and $B$ belong to the subcategory in the statement.
Apply the preceding argument to each shifted triangle: its two inputs
are $F(\Sigma^jB)$ and $F(\Sigma^{j+1}A)$, both in $\mathcal D_0$.
It follows that $F(\Sigma^jC)\in\mathcal D_0$ for every $j$.
Rotating triangles gives the other two-out-of-three implications.
\end{proof}

Even under this weaker hypothesis, the equivalence preserves addition
on Hom sets: addition is determined by biproduct diagonals and folds.
Consequently an identification of maps out of the sphere obtained from
such a comparison is an identification of abelian groups, as well as
of their underlying sets.
